\documentclass[10pt,a4paper]{article}

\usepackage[utf8]{inputenc}
\usepackage[T1]{fontenc}
\usepackage{babel}
\babelprovide[main,import]{portuguese}
\usepackage[margin=2cm]{geometry}
\usepackage{array,tabularx,booktabs}
\usepackage{titlesec}
\usepackage{tikz}
\usetikzlibrary{positioning,arrows.meta,fit}

\setlength{\parindent}{0pt}
\setlength{\parskip}{3pt}
\setlength{\emergencystretch}{2em}
\renewcommand{\arraystretch}{1.08}
\newcolumntype{Y}{>{\raggedright\arraybackslash}X}
\newcolumntype{P}[1]{>{\raggedright\arraybackslash}p{#1}}
\titlespacing*{\section}{0pt}{9pt}{4pt}
\pagestyle{empty}

\title{Proposta de Projeto: \textit{Egg Plotter}}
\author{Guilherme Rodrigues dos Santos, Laura Armiliato Sangalli, Lucas Tarcisio Morais Pies}
\date{Setembro de 2026}

\begin{document}

%  Introdução
% Definição: hardware e software
% Cronograma
%  Análise de custos e riscos
%
%
%
%


\maketitle

\section{Introdução}

Desenhar sobre uma superfície plana é um problema resolvido há décadas por
impressoras e \textit{plotters}. Desenhar sobre um ovo não é: a superfície
é curva, não pode ser planificada sem distorção e muda de diâmetro ao longo
do comprimento. O \textit{Egg Plotter} trata essa dificuldade com dois
movimentos de rotação: o objeto gira em torno do próprio eixo enquanto um
braço inclina a caneta de uma extremidade à outra, de modo que cada ponto
da superfície é alcançado pela combinação dos dois ângulos. Um servomotor
levanta e abaixa a caneta entre os traços.

Máquinas desse tipo existem como projetos abertos, e a equipe já montou
uma a partir do EggBot/SphereBot e do firmware EggDuino. Nelas, o desenho
é preparado em um computador, com programas como o Inkscape, e enviado ao
Arduino por cabo USB. Em uma feira, esse arranjo reduz o público a
espectador, já que apenas quem opera o computador consegue desenhar.

A proposta é evoluir essa solução com um ESP32, que passará a criar uma
rede Wi-Fi local e a hospedar uma interface web, pela qual qualquer
visitante poderá desenhar ou importar uma figura no navegador de um
celular ou computador, conferir a prévia e enviá-la ao equipamento, sem
instalar aplicativos e sem depender de internet. A contribuição da equipe
concentra-se nessa camada: interface, processamento do desenho e
comunicação sem fio.

O sistema se destina a atividades de divulgação científica, como a Feira
de Profissões da UTFPR e a Inventum. Nelas, o visitante acompanha o
processo completo da execução do desenho, o que
torna o equipamento uma forma concreta de apresentar a integração entre
eletrônica, programação e controle de movimento.

As seções seguintes apresentam a definição do hardware e do software, o
cronograma de desenvolvimento e a análise de custos e riscos.


\section{Definição: hardware \& software}

O \textit{Egg Plotter} envolverá conceitos relacionados à programação web, sistemas embarcados, vetorização de imagens e acionamento de atuadores. Nesse sentido, conceitos de hardware e software serão aplicados e apresentarão interconexão ao longo do projeto. A seguir, a descrição das funções executadas por essas duas áreas será apresentada.

\subsection{Hardware}

Serão reaproveitados a estrutura mecânica, o Arduino e o acionamento dos
eixos existentes. A nova etapa concentra-se na integração do ESP32, no
controle da caneta e na validação do conjunto. A Tabela~\ref{tab:hardware}
apresenta os componentes previstos; os modelos e conexões serão conferidos
na primeira sprint.





\begin{table}[htbp]
\centering
\caption{Componentes de hardware.}
\label{tab:hardware}
\small
\begin{tabularx}{\linewidth}{@{}l c X@{}}
\toprule
\textbf{Componente} & \textbf{Qtd.} & \textbf{Função} \\
\midrule
Arduino UNO R3 & 1 & Executa o firmware EggDuino e gera os pulsos de passo. \\
CNC Shield v3 & 1 & Distribui sinais e alimentação aos drivers; micropasso configurado em 1/16 por jumpers. \\
Driver A4988 (ou DRV8825) & 2 & Acionamento dos motores de passo dos eixos X (rotação) e Y (braço). \\
Motor de passo NEMA~17 & 2 & Rotação do ovo e deslocamento angular do braço da caneta (3200 passos/volta com micropasso). \\
Servo SG90 & 1 & Levanta e abaixa a caneta (pino D4). \\
ESP32 DevKit & 1 & Rede Wi-Fi, servidor web, geração de comandos e controle da plotagem. \\
Fonte 12\,V & 1 & Alimentação dos motores via CNC Shield. \\
Divisor resistivo (1\,k$\Omega$ / 2\,k$\Omega$) & 1 & Adequa o sinal TX do Arduino (5\,V) à entrada RX do ESP32 (3,3\,V). \\
Estrutura mecânica & 1 & Peças impressas em 3D, rolamentos, eixos, mola de pressão da caneta. \\
Caneta marcadora & 1 & Ferramenta de desenho. \\
\bottomrule
\end{tabularx}
\end{table}


\subsection{Software}

O software está organizado em três camadas, uma por bloco de hardware:

\begin{enumerate}
  \item \textbf{Cliente web (HTML + CSS + Javascript)}: tela de desenho em \textit{canvas} e envio de imagem; vetorização de imagens enviadas; simplificação dos traços pelo algoritmo de Ramer--Douglas--Peucker; ordenação dos traços para reduzir deslocamentos com a caneta levantada; pré-visualização do resultado na superfície desenrolada do ovo; acompanhamento do progresso.
 \item \textbf{Servidor no ESP32 (C++/Arduino)}: ponto de acesso Wi-Fi (\texttt{WiFi.softAP}); servidor HTTP assíncrono com os arquivos da interface armazenados em LittleFS; API para receber trajetórias e comandar início, pausa e cancelamento; conversão das coordenadas em passos dos eixos (área útil de $3200\times800$ passos) e em comandos EBB (\texttt{SM} para movimento, \texttt{SP} para a caneta, \texttt{EM} para habilitar motores); envio controlado por fila, verificando se a resposta \texttt{OK} do firmware indica recebimento ou fim do movimento.
  \item \textbf{Firmware EggDuino (C++)}: \textit{fork} \texttt{plex3r/\allowbreak EggDuino}, com as bibliotecas AccelStepper e SerialCommand; recebe os comandos pela serial a 9600\,baud e executa os movimentos.
\end{enumerate}

Como ferramenta de desenvolvimento e validação, o repositório mantém o \textit{pré-plotter} em Python (\texttt{preplotter/}), que envia padrões geométricos (senoide, círculo, espiral, estrela) diretamente ao Arduino por USB a partir de um computador. Ele serve de referência para validar a comunicação e a calibração antes que a mesma lógica seja portada para o ESP32.

% ------------------------------------------------------------
\section{Cronograma}

Propõem-se quatro sprints de duas semanas, totalizando oito semanas. Cada ciclo terá uma
entrega demonstrável, permitindo verificar o funcionamento antes de ampliar
o escopo. A montagem existente e a recepção serial são o ponto de partida,
e não atividades a refazer integralmente.

\begin{center}
\small
\begin{tabularx}{\linewidth}{@{}P{1.1cm}P{1.6cm}Y@{}}
\toprule
\textbf{Sprint} & \textbf{Período} & \textbf{Atividades e critério de entrega}\\
\midrule
1 & Semanas 1--2 & Conferir hardware e protocolo, organizar o firmware do ESP32, realizar a ligação UART e calibrar os eixos. \textbf{Entrega:} ESP32 reproduz movimentos de referência nos dois eixos, com sequência e limites documentados.\\
2 & Semanas 3--4 & Criar o Wi-Fi local, hospedar a interface, implementar comunicação HTTP, transferência de trabalhos e controle manual. \textbf{Entrega:} celular ou computador inicia um trabalho de teste pelo navegador e acompanha o estado.\\
3 & Semanas 5--6 & Implementar desenho livre, importação SVG, simplificação, prévia e conversão das trajetórias em movimentos. \textbf{Entrega:} um desenho contínuo é enviado pela interface e reproduzido no objeto.\\
4 & Semanas 7--8 & Finalizar controle da caneta, importação básica de imagens, pausa e cancelamento. Testar falhas, corrigir problemas e documentar. \textbf{Entrega:} desenho com traços separados, demonstração integrada e guia de operação.\\
\bottomrule
\end{tabularx}
\end{center}

As tarefas serão acompanhadas no Jira, com responsável, dependências e
critérios de aceite. Ao final de cada sprint, a equipe registrará os testes
e apresentará o resultado. As estimativas serão revistas conforme a
capacidade da equipe; recursos de importação terão seu nível de sofisticação
ajustado antes de comprometer a integração básica. A disponibilidade do
mecanismo da caneta deverá ser verificada já na primeira sprint.


\section{Análise de custos e riscos}
A seguir, serão apresentados os componentes do sistema que implicam custo na execução do projeto, bem como os riscos relacionados à possíveis falhas de hardware, planejamento ou execução.
\subsection{Custos}

A Tabela~\ref{tab:custos} apresenta a estimativa de custo dos componentes. Os valores são referências de mercado nacional e devem ser confirmados com orçamento; parte dos itens já foi adquirida.

\begin{table}[htbp]
\centering
\caption{Estimativa de custos.}
\label{tab:custos}
\small
\begin{tabular}{@{}l c r r l@{}}
\toprule
\textbf{Item} & \textbf{Qtd.} & \textbf{Unit. (R\$)} & \textbf{Total (R\$)} & \textbf{Situação} \\
\midrule
Arduino UNO R3                          & 1 & 180,00 & 180,00 & adquirido \\
CNC Shield v3                           & 1 &  30,00 &  30,00 & adquirido \\
Driver A4988                            & 2 &  15,00 &  30,00 & adquirido \\
Motor de passo NEMA~17                  & 2 &  70,00 & 140,00 & adquirido \\
Servo SG90                              & 1 &  20,00 &  20,00 & adquirido \\
ESP32 DevKit                            & 1 &  50,00 &  50,00 & a integrar \\
Fonte 12\,V                             & 1 &  45,00 &  45,00 & adquirido \\
Peças impressas em 3D (filamento PLA)   & 1 &  40,00 &  40,00 & adquirido \\
Rolamentos, parafusos, mola, eixos      & 1 &  40,00 &  40,00 & adquirido \\
Resistores, jumpers, conectores         & 1 &  10,00 &  10,00 & a adquirir \\
Canetas marcadoras                      & 1 &  15,00 &  15,00 & a adquirir \\
Ovos para teste                         & 1 &  20,00 &  20,00 & a adquirir \\
\midrule
\multicolumn{3}{@{}l}{\textbf{Subtotal}}          & \textbf{620,00} & \\
\multicolumn{3}{@{}l}{Reserva de contingência (15\,\%)} & 93,00 & \\
\multicolumn{3}{@{}l}{\textbf{Total estimado}}    & \textbf{713,00} & \\
\bottomrule
\end{tabular}
\end{table}

O software utilizado é integralmente livre ou gratuito (Arduino IDE, bibliotecas AccelStepper e SerialCommand, firmware EggDuino, Python, navegador web), sem custo de licença. Caso o plano de contingência R1 (Tabela~\ref{tab:riscos}) seja acionado, soma-se o custo de um Raspberry Pi.

\subsection{Riscos}

Os riscos foram classificados por probabilidade (P) e impacto (I) em três níveis: baixo (B), médio (M) e alto (A).

\begin{table}[htbp]
\centering
\caption{Análise de riscos.}
\label{tab:riscos}
\small
\begin{tabularx}{\linewidth}{@{}l X c c X@{}}
\toprule
\textbf{ID} & \textbf{Risco} & \textbf{P} & \textbf{I} & \textbf{Mitigação / contingência} \\
\midrule
R1 & Memória e processamento do ESP32 insuficientes para o processamento das imagens. & M & A & Vetorização e simplificação executadas no navegador; ESP32 recebe apenas trajetórias. Contingência: substituir o ESP32 por um Raspberry Pi, reaproveitando a interface web. \\
R2 & Incompatibilidade de tensão (5\,V/3,3\,V) na UART danificar o ESP32. & M & A & Divisor resistivo ou conversor de nível no sentido Arduino$\rightarrow$ESP32; conferir com multímetro antes de ligar. \\
R3 & Conflito da UART D0/D1 com a gravação de firmware e o monitor serial. & A & B & Desconectar o ESP32 durante a gravação; firmware definitivo gravado antes da integração final. \\
R4 & Reinicialização do Arduino (autoreset) ao abrir a serial, perdendo a posição da caneta. & M & M & Desabilitar o autoreset (capacitor de 10\,$\mu$F entre RESET e GND) apenas após o firmware definitivo, pois impede novas gravações. \\
R5 & Eixos trocados na geração dos comandos \texttt{SM}. & M & M & Teste com movimento em um único eixo antes do primeiro desenho; parâmetro de configuração para inverter a ordem. \\
R6 & Falha na gravação do Arduino com o CNC Shield encaixado. & A & B & Procedimento definido: remover o shield antes de cada gravação. \\
R7 & Variação de forma do ovo, escorregamento ou pressão inadequada da caneta. & A & M & Calibração da área útil, ajuste da mola e testes em ovos cozidos ou de plástico. \\
R8 & Aquecimento dos drivers e perda de passos. & B & M & Ajuste da corrente (Vref) dos drivers e uso de dissipadores. \\
R9 & Atraso na integração entre os módulos. & M & A & Organização de entregas incrementais; pré-plotter em Python como alternativa de demonstração. \\
\bottomrule
\end{tabularx}
\end{table}
\end{document}
